\documentclass[11pt,a4paper]{article}
\usepackage[margin=2.2cm]{geometry}
\usepackage{newtxtext,newtxmath}
\usepackage[T1]{fontenc}
\usepackage{microtype}
\usepackage{booktabs,tabularx,array}
\usepackage{enumitem}
\usepackage{xcolor}
\usepackage{titlesec}
\usepackage[hidelinks]{hyperref}

\definecolor{ink}{HTML}{083D5E}
\definecolor{accent}{HTML}{2689CA}
\titleformat{\section}{\large\bfseries\color{ink}}{\thesection}{0.6em}{}
\titleformat{\subsection}{\normalsize\bfseries\color{ink}}{\thesubsection}{0.5em}{}
\titlespacing*{\section}{0pt}{1.1em}{0.4em}
\titlespacing*{\subsection}{0pt}{0.8em}{0.3em}
\setlist{itemsep=1pt,topsep=3pt,leftmargin=1.3em}
\setlength{\parindent}{0pt}
\setlength{\parskip}{0.45em}
\newcolumntype{L}{>{\raggedright\arraybackslash}X}
\renewcommand{\arraystretch}{1.15}
\newcommand{\say}[1]{\par\smallskip\noindent\colorbox{accent!10}{\parbox{\dimexpr\linewidth-2\fboxsep}{\small\itshape #1}}\par}

\begin{document}

\begin{center}
{\LARGE\bfseries\color{ink} Adaptive Task Allocation for Heterogeneous Bio-Inspired Agents}\\[0.4em]
{\large The whole project, explained}\\[0.3em]
{\small Group GP105 \,·\, Nikshay Yadav, Madan Raj Murugesan, Aryan Varshney \,·\, Supervisor: Prof.\ Shailender Kumar}\\[0.2em]
{\small Panel preparation, 17 November 2026}
\end{center}

\section{The idea in one paragraph}
Many AI systems now use several language-model agents working together: one plans, others solve parts,
others check or vote. Almost all of them pick \emph{one} way of working together (a debate, a pipeline, a
manager with workers) and use it for the whole task. Real tasks are mixed: one part is planning, one is
maths, one is routine. Our question: can a coordinator that \textbf{picks a different way of working for
each part of the task} do better than one strong agent working alone, \textbf{at exactly the same
compute}? We took the ``ways of working'' from nature: ant colonies, honeybee swarms and the octopus.

\section{Why this is worth studying}
\begin{itemize}
  \item \textbf{Kim et al.\ (2026):} same model, same compute; switching from one agent to a multi-agent
        design changed results by anywhere from $+80.8\%$ to $-70\%$ depending on the task. The
        coordination choice matters a lot, and no single choice is always best.
  \item \textbf{Cemri et al.\ (2025):} many failures of popular multi-agent frameworks come from
        coordination, not from the model being weak.
  \item \textbf{Wang et al.\ (2024):} one agent with a strong prompt comes close to the best
        multi-agent discussions, so ``more agents'' is not automatically better.
  \item \textbf{Compute confound:} papers often give the multi-agent system far more compute than the
        single agent, so the ``win'' is just extra spending.
  \item \textbf{Metaphor criticism:} many nature-inspired algorithms are ordinary techniques with an
        animal name on top (Camacho-Villal\'on et al.), so we test every mechanism against a version
        with the biological part removed.
\end{itemize}

\section{The gaps we address}
Existing adaptive systems choose an agent graph (GPTSwarm), a communication topology (G-Designer), a
workflow (AFlow), a model (RouteLLM), collaboration modes and roles (MasRouter), or which agent gets a
subtask with a LinUCB bandit (Symphony-Coord). AMRO-S and SwarmSys apply ant-colony ideas to LLM agents.
\begin{itemize}
  \item \textbf{G1:} nobody chooses between whole \emph{coordination mechanisms}, \emph{per subtask}.
  \item \textbf{G2:} comparisons are often made at unequal compute.
  \item \textbf{G3:} bio-inspired methods are rarely tested against controls without their biological part.
\end{itemize}

\section{Research questions and status after Phase I}
\begin{tabularx}{\linewidth}{@{}p{1.0cm}p{6.2cm}L@{}}
\toprule
\textbf{RQ} & \textbf{Question} & \textbf{Status} \\
\midrule
RQ1 & Do the mechanisms have different strengths? & Yes: the single agent wins planning, the bee wins maths \\
RQ2 & Does choosing per subtask beat any fixed choice at the same budget? & Shown with a simple by-type rule; learned version in Phase II \\
RQ3 & What does the biological part add over a plain version? & The bee's quorum is a reliable confidence signal; recruitment only adds cost \\
RQ4 & Can the system recover from failures? & Reallocation built; failure-injection tests in Phase II \\
RQ5 & When does coordination stop paying for itself? & Voting costs about $5\times$ per correct answer \\
\bottomrule
\end{tabularx}

\section{The mechanisms}
All of them are fixed coordination rules around language-model calls. The rule decides who works, how
many, when to stop and whose answer wins; the language model does the solving.
\begin{itemize}
  \item \textbf{Single agent (main baseline):} one conversation. It answers, gets the check's feedback
        (e.g.\ ``plan invalid at step 7 because\ldots''), fixes the answer, and repeats.
  \item \textbf{Direct:} one model call, no retries. The cheapest possible system, used as a floor.
  \item \textbf{Ant colony:} builds a solution step by step, choosing each step with a probability based on
        ``pheromone'' weights. Steps in good solutions get more pheromone, and all pheromone slowly
        evaporates (classic ant colony optimisation, Dorigo). \emph{Controls:} no deposit, no evaporation,
        uniform weights. A ``macro ant'' chooses chunks of plan instead of single actions.
  \item \textbf{Honeybee (quorum bee):} based on how bee swarms pick a nest (Seeley). Five scouts solve the
        problem independently. As soon as \textbf{3 of 5 agree} (the quorum), or an answer passes a check that
        proves it, it stops. Without a quorum: \emph{cross-inhibition} (committed scouts may be released
        when a rival exists, probability 0.2) and \emph{recruitment} (uncommitted scouts look at another
        scout's answer, chosen by support, and may adopt it). \emph{Controls:} no stop signal, and a plain
        \textbf{majority vote over five answers} (the key control).
  \item \textbf{Octopus:} based on the octopus arm finishing a reach on its own once the brain sends the
        command (Sumbre et al.). The coordinator hands a bounded job to one ``arm'', which works locally
        with check feedback and reports back only the result. \emph{Control:} per-step reporting.
\end{itemize}

\section{The framework: \texttt{bioalloc}}
About 9,400 lines of Python and 5,800 lines of tests (372 tests), checked with ruff, strict mypy and
pytest on GitHub Actions. Every design decision (D1--D33) and every review or result (R1--R29) is
recorded in \texttt{docs/decisions.md}.

\subsection{How a task flows through it}
\begin{enumerate}
  \item A task arrives as a \textbf{graph of subtasks}, decomposed in advance and stored, so every system
        gets exactly the same breakdown.
  \item The \textbf{coordinator} takes the subtasks whose dependencies are done and asks the
        \textbf{allocation policy} which mechanism should take each (constant, random, rule table, LinUCB).
  \item The \textbf{executor} (ant, bee, octopus, direct, \ldots) runs the subtask.
  \item \textbf{Runtime checks} test the result while the system is running. For planning, our STRIPS
        simulator runs the plan; for maths, the check can only confirm the answer is a single number.
        If the result passes, it is accepted; if it fails, it goes back to the coordinator with the reason and
        is reallocated, up to a retry limit that is the same for every system.
  \item The \textbf{aggregator} combines the subtask results into the final answer.
  \item A separate \textbf{grader}, the only module holding the correct answers, scores the result
        \emph{after} the run. A test enforces that nothing else imports it, so no system can peek.
\end{enumerate}

\subsection{The token budget}
The budget belongs to the \emph{task}, not to the mechanism: every system gets the same cap for the same task.

\begin{tabularx}{\linewidth}{@{}Lrr@{}}
\toprule
\textbf{Limit per task} & \textbf{Planning (Logistics)} & \textbf{Maths (AIME)} \\
\midrule
Tokens & 21,000 & 60,000 \\
Model calls & 60 & 30 \\
Wall time & 10 min & 15 min \\
Tool calls & 40 & 40 \\
\bottomrule
\end{tabularx}

Every call goes through a \textbf{metered client}: it reserves a worst-case token cost before sending and
\textbf{refuses the call} if that does not fit. Subtasks draw from the same task pot. The cap is a
ceiling, not an allowance: systems can use less, so we also report the cost per correct answer.

\subsection{Logging}
Every run writes a trace of every call and a manifest with the config, code version and data-split hash.
Any number in the report can be traced to \texttt{runs/*/results.jsonl}.

\section{Data}
Each suite is stored as a \emph{public} file (what systems see) and a \emph{private} file (answers, read
only by the grader), split once into frozen train, tune and test sets.

\begin{tabularx}{\linewidth}{@{}p{3.6cm}p{1.3cm}Lp{3.0cm}@{}}
\toprule
\textbf{Suite} & \textbf{Size} & \textbf{Content} & \textbf{Judged by} \\
\midrule
GSM8K + our tool suite & 150+150 & Word problems; multi-tool tasks & Number / exact match \\
PlanBench & 150 & Blocksworld and Logistics planning (PDDL) & Our STRIPS simulator \\
Generated Logistics (ours) & 180 & 3--4 cities, 5/7/9 packages & Simulator \\
AIME 2020--2025 & 178 & Competition maths; 30 tuning, \textbf{148 held-out test} & Exact match after the run \\
\bottomrule
\end{tabularx}

\section{How we ran experiments}
\begin{itemize}
  \item \textbf{Pre-registration:} before each experiment we wrote down in \texttt{decisions.md} what we
        would measure, what counts as a win and what we would conclude, and only then ran it. This stops us
        from fitting the story to the results.
  \item \textbf{Stages:} smoke test (cheap check that nothing is broken) $\rightarrow$ calibration (find tasks
        where the single agent is neither always right nor always wrong) $\rightarrow$ comparison at equal
        budget $\rightarrow$ tuning with an advancement rule fixed in advance $\rightarrow$ held-out test only
        for variants that advance.
  \item \textbf{Statistics:} differences use a \textbf{paired bootstrap with 95\% confidence intervals}.
        Both systems ran the same problems; we resample the problems thousands of times with replacement,
        recompute the gap each time, and keep the middle 95\% of gaps. If the interval does not include
        zero, the gap is very unlikely to come from which problems we happened to test.
  \item \textbf{Models:} Claude Sonnet 4.6 on Amazon Bedrock (most of Phase I) and GPT-6 Luna on OpenAI
        (about $30\times$ cheaper; used to replicate the key results and for Phase II). Both at temperature
        1.0, low reasoning effort.
\end{itemize}

\section{What happened, in order}

\subsection*{Step 1: Building the framework (October)}
Everything was first built offline against a \textbf{scripted fake model} so the logic could be tested for
free. Then came the budget ledger, checks, executors, task suites and splits. After every milestone an
independent code review (by another AI model) was run; findings were fixed and logged.

\subsection*{Step 2: Choosing a backbone}
The plan started with GPT-OSS-120B on Groq (D3); the pilot moved to Claude Sonnet 4.6 on Bedrock (D23).

\subsection*{Step 3: Calibration (R12)}
At a generous 120,000-token cap the single agent solved \emph{every} calibration task: normal maths and
easy planning were too easy for any mechanism to beat it. On Logistics it solved 10 of 15 at 21,000 tokens
and 15 of 15 at 50,000. So we chose \textbf{Logistics at 21,000 tokens} (about twice the single agent's
average spend), where there is room to improve.

\subsection*{Step 4: Logistics comparison on Sonnet (R13)}
Single agent 73\%, octopus 67\%, bee 60\%, direct 57\%, ant 0\%. No mechanism beat the single agent. The
traces showed why:
\begin{itemize}
  \item The single agent reads the simulator's error message and fixes its plan.
  \item The bee's scouts each made one plan and never saw why it was rejected; recruitment never even ran.
  \item The ant spent one model call per action. Plans need 14--45 actions and the budget buys about 5--7
        calls, so it never finished a plan.
\end{itemize}

\subsection*{Step 5: Trying to fix the mechanisms for planning}
\begin{itemize}
  \item \textbf{Repair scouts (D28, R21):} scouts that fix their own rejected plans using the simulator's
        feedback. Three variants ran on 30 tuning instances of our generated suite. The rule fixed in advance:
        advance only if a variant beats the single agent, or matches it with 20\% fewer tokens. The best came
        within one success (23/60 vs 24/60) but used 4\% \emph{more} tokens, so nothing advanced. The 120
        held-out generated instances remain unseen, reserved for Phase II.
  \item \textbf{Macro ant (D30, R22):} 0--1 of 12 solved. The model mostly proposed illegal plan chunks, so
        pheromone had nothing to choose between. Lesson: ant-style search fits choices between several
        cheap, checkable options, not long plans.
\end{itemize}

\subsection*{Step 6: Where the bee should win: AIME maths (D31)}
The bee's real strength is a quorum of independent opinions, which only helps when \textbf{nothing can
check the answer}. AIME is exactly that: one integer answer, no checker. The \textbf{first test was stopped
(R23)}: both voters scored below the single agent, which made no sense. The traces showed a fairness bug:
each scout was limited to 4 model replies, so scouts that used the calculator a few times were cut off with
no answer, while the single agent had 12 steps. We fixed it (12 for scouts too), disclosed that 52 test
problems had been seen in the stopped run, and re-ran.

\subsection*{Step 7: Switching to Luna (D32)}
Sonnet was expensive (US\$3 / \$15 per million input/output tokens) and its daily quota kept killing long
runs. GPT-6 Luna costs US\$0.10 / \$0.50. Sonnet results stay Sonnet results; models are never mixed inside
one experiment.

\subsection*{Step 8: AIME held-out test on Luna (R26), 148 problems}
\begin{tabularx}{\linewidth}{@{}Lrrr@{}}
\toprule
\textbf{System} & \textbf{Solved} & \textbf{Tokens/problem} & \textbf{Gap vs single agent} \\
\midrule
Single agent & 70.9\% & 1,112 & -- \\
Majority of five & 84.5\% & 6,004 & $+13.5$ (CI $+5.4$ to $+21.6$) \\
\textbf{Quorum bee} & \textbf{86.5\%} & 6,441 & $\mathbf{+15.5}$ (CI $+8.1$ to $+23.0$) \\
\bottomrule
\end{tabularx}

Both voters clearly beat the single agent. Bee vs majority is $+2.0$ points (CI $-2.0$ to $+6.8$), i.e.\
noise, so \textbf{most of the gain is self-consistency} (voting), at about $5\times$ the cost per correct
answer (still well under a cent). On Sonnet (stopped at 141 of 148 problems) the order was the same.

\subsection*{Step 9: Logistics replicated on Luna (R27)}
15 problems $\times$ 2 runs = 30 runs per system. Single agent 27/30 (90\%), octopus 87\%, bee 14/30 (47\%),
direct 47\%, ant 0\%: \textbf{the same order as on Sonnet}. Single agent minus bee: $+43$ points
(13/30), \textbf{95\% CI $+20$ to $+67$}. The interval is wide because 15 problems is a small set, but it
stays well above zero. Caveat we report ourselves: we had pre-set a ``near-ceiling'' flag at 90\%, and the
single agent hit exactly 90\%; the bee, direct and ant are far below it, so the problems still separate the systems.

\subsection*{Step 10: What the bee adds beyond voting (R28, analysed after the test)}
\begin{itemize}
  \item \textbf{Quorum stop:} on \textbf{120 of 148} problems three scouts agreed early. The bee stopped and
        was \textbf{right 97\%} of the time, using $0.69\times$ a full vote's tokens. The quorum is a cheap,
        reliable ``I am confident'' signal.
  \item \textbf{Recruitment:} on the other 28 problems recruitment solved 43\% (a plain majority 50\%) at
        $1.87\times$ the tokens. Recruitment adds cost, not accuracy.
  \item \textbf{Phase II hypothesis (untested):} stop at a quorum, otherwise take a plain vote, using about
        $0.75\times$ a full vote's tokens.
\end{itemize}

\subsection*{Step 11: Choosing by task type (R29)}
\begin{tabularx}{\linewidth}{@{}Lrrr@{}}
\toprule
\textbf{Policy (GPT-6 Luna)} & \textbf{Planning} & \textbf{Maths} & \textbf{Mean} \\
\midrule
Always single agent & 90.0\% & 70.9\% & 80.5\% \\
Always bee & 46.7\% & 86.5\% & 66.6\% \\
\textbf{Choose by task type} & \textbf{90.0\%} & \textbf{86.5\%} & \textbf{88.2\%} \\
\bottomrule
\end{tabularx}

This is the project's central claim in miniature: the best mechanism flips with the kind of task, so
choosing per subtask beats any fixed choice. Caveat: measured on the same runs (in-sample), so it is a
demonstration, not a test. Phase II learns the choice and tests it on unseen tasks.

\textbf{Cost of Phase I:} about US\$125 on Sonnet and US\$1.60 on Luna. A full study would cost
US\$1,000--1,900 at Sonnet prices and under US\$100 on Luna.

\section{The three takeaways}
\begin{enumerate}
  \item \textbf{When a task has a checker} (planning), a single agent that learns from feedback is very hard
        to beat; none of our mechanisms did.
  \item \textbf{When nothing checks the answer} (maths), voting beats the single agent by 13--16 points, and
        the bee's quorum tells you when to trust the answer (97\% right when reached).
  \item \textbf{The best mechanism depends on the task type}, which is exactly why per-subtask allocation matters.
\end{enumerate}

\section{Limitations (say these before the panel does)}
\begin{itemize}
  \item Each task type is one suite (Logistics, AIME).
  \item The by-type result is in-sample.
  \item The quorum analysis was made after the test.
  \item 15 PlanBench problems detect only large effects, so the intervals are wide.
  \item The ant never got a fair use case in Phase I.
\end{itemize}

\section{Phase II (December--May)}
\begin{itemize}
  \item Build \textbf{mixed task graphs}: planning, calculation and tool use in one task.
  \item \textbf{Learn the allocation policy}: a rule table first, then a \textbf{LinUCB contextual bandit}
        that learns which mechanism to pick from subtask features and past outcomes.
  \item Test the \textbf{quorum-stop-then-vote} idea.
  \item Add two more mechanisms, and \textbf{failure injection} (RQ4): break agents or tools and check that
        reallocation recovers.
  \item Run the full study with ablations on held-out tasks; write the thesis and a paper.
\end{itemize}

\section{Likely panel questions}

\subsection*{``What exactly is a runtime check?''}
An automatic test on a result while the system is still running, using only what would exist in real use:
the plan simulator for planning, a format check (``is it a single number?'') for maths. Pass: accepted.
Fail: sent back with the reason and reallocated. Runtime checks never see the correct answers; only the
grader does, after the run.

\subsection*{``Which comparison is fair, and how?''}
Each bio-inspired mechanism against the plain versions: mainly the single agent, plus the majority vote and
the direct call. The only difference allowed is how the agents work together: same model, prompt, checks
and token cap for every system, enforced in code, and no system ever sees the answers.

\subsection*{``Is there a separate token cap for the bee, the octopus, \ldots?''}
\say{Every system gets the same budget for each task, 21,000 tokens for a planning task and 60,000 for a
maths problem. It's a ceiling: a system can use less, but every call is checked before it's sent, and any
call that would go over is refused. So no mechanism can win just by spending more.}

\subsection*{``Aren't these just deterministic algorithms you call agents?''}
\say{That's fair for the coordination layer. The ant, bee and octopus rules are deliberately written as fixed
algorithms, so we can test whether the mechanism itself helps. But the agents inside them are language-model
calls that each reason independently and give different answers every run, and that's what the rules
coordinate. It's the same structure as debate or manager--worker systems in the literature. And the
adaptive part of the project is choosing which mechanism to use for each subtask, which Phase II learns from data.}

If pushed on ``why bio-inspired then?'': the biology is the source of the coordination rules (quorum
sensing, pheromone reinforcement, decentralised arms); we do not claim the agents are bees.

\subsection*{``What does $+80.8\%$ to $-70\%$ mean?''}
\say{One study tested the same AI working alone versus as a team. On some tasks the team did 80 percent better;
on others it did 70 percent worse. So teamwork isn't always better; it depends on the task. That's the idea
behind our project: pick the right way of working for each part of the task.}

\subsection*{``Where do the 43 points and the 20--67 interval come from?''}
\say{On Luna the single agent made a valid plan in 27 of 30 runs and the bee in 14, a 43-point gap. We resampled
the problems thousands of times to see how much that gap could move by chance; it stayed between 20 and 67
points, never near zero, so the difference is real, not an accident of which problems we picked.}

\vfill
\small Code and results: \texttt{github.com/Nikshay1/BTP} \,·\, Report: \texttt{report/BTP-Phase1-Report-GP105.pdf}
\,·\, Decision log: \texttt{docs/decisions.md}

\end{document}
